\section{Discussion}\label{sec:discussion}

\subsection{Why the VPP--guidance interface matters}

The current LOS-rate guidance law is driven by the virtual point selected by the high-level policy. In the default implementation, target information reaches the low-level controller mainly through that geometric point; target velocity is used for speed regulation, but target maneuver state is not directly used to shape the lateral or normal-acceleration command. This design is interpretable, but it creates a narrow interface. If the high-level policy chooses a virtual point that is dynamically awkward, the low-level law has limited ability to recover from that choice.

This explains why the most useful interpretation of the current evidence is not ``VPP is better'' or ``prediction fails.'' A predictor may estimate future target motion correctly, yet its benefit can be lost if the only downstream channel is an offset point that the guidance law cannot track safely. Conversely, pure PN without VPP can recover selected tail-chase candidates, which means some failures cannot be assigned to physical infeasibility. The remaining hypothesis is an interface-level one: the mapping from predictor output to virtual point and from virtual point to LOS-rate commands may be the active bottleneck under the tested reward and geometry envelope. The completed target-state-direct ablation does not rescue the disadvantage case, which narrows the diagnosis but does not close it; a redesigned guidance interface or PN-gated mode-switch experiment is still needed.

\subsection{Negative-result value}

The main value of this revision is to make the negative result publishable rather than accidental. Under nominal simple-backend cases, the zero-offset guidance chain can match or exceed VPP variants. Under corrected JSBSim regression, no-prediction and gain-only variants share the same 50\% success rate. Under the new formal external-baseline and interface-ablation matrices, pure pursuit, pure PN, APN1971, full VPP, zero-offset No-VPP, and target-state-direct guidance all share the same success/failure pattern on the frozen feasible-geometry suite: complete success in favorable, neutral, and challenging scenarios, and complete out-of-bounds failure in the disadvantage scenario. Taken together, these observations argue against a simple algorithm-ranking story. They point instead to a fair-comparison problem: the paper must compare controllers under the same observations, seeds, simulator, metrics, and training budget before claiming where the bottleneck lies.

\subsection{Statistical interpretation}

Deep-RL seed sensitivity is not a nuisance detail in this study; it is part of the result. Episode-level exact tests can show that a terminal-outcome difference is unlikely under pooled counts, but they do not replace seed-level reproducibility. The final Defence Technology submission should therefore report both levels: episode-level Fisher or McNemar tests for binary outcomes, and seed-level confidence intervals or paired tests for training-run stability. Claims based on fewer than 10 training seeds should remain exploratory unless the effect is reproduced across independent checkpoints.

\subsection{Limitations}

Four limitations bound the current manuscript.

\begin{enumerate}
    \item \textbf{Incomplete final matrix:} the repository contains useful 30-evaluation-seed simple-backend data, a corrected 10-evaluation-seed JSBSim sanity run, and new 3-seed formal external-baseline/interface-ablation matrices, but not the full 10-training-seed JSBSim matrix for VPP, No-VPP, E2E, PN, APN, and TPP-style baselines.
    \item \textbf{Incomplete trained ablations:} no-regret, no-gain-observation, and no-safety-penalty checkpoints are still absent. The current ablation evidence covers zero-offset No-VPP and target-state-direct guidance, not the full mechanism-removal set.
    \item \textbf{Scenario coverage:} current scenarios cover favorable, neutral, disadvantage, and challenging geometries, but defence-context claims require head-on, offset crossing, stern conversion, maneuvering tail-chase, wind/noise/delay robustness, and parameter-sensitivity sweeps.
    \item \textbf{Reward dependence:} timeout and crash trade-offs depend strongly on terminal penalties and safety weights in \texttt{config/reward.yaml}. Different mission preferences may change the interpretation of conservative VPP behavior.
\end{enumerate}

\subsection{Future work}

The next technical step is not another policy variant. It is an interface redesign experiment. Candidate directions include direct target-state conditioning in the guidance law, PN-gated mode switching with latch logic, curriculum learning over stern-conversion geometry, and a VPP action space constrained by dynamic feasibility. The robustness script added in this revision also makes it possible to quantify wind, sensor noise, actuator delay, and guidance-gain sensitivity under a common metric schema. These experiments are prerequisites for moving from a credible negative result to a deployable controller claim.
